\section{Illustrated Notation Atlas}
\label{sec:interface-notation-atlas}

This appendix supplies examples for the two notation-heavy interfaces in
Sections~\ref{sec:reader-framework}.  Each panel is a separate TikZ source so
that it can be retained or removed independently.  The diagrams explain the
definitions and hypotheses; they are not substitutes for the inequalities in
Propositions~\ref{prop:body-transfer-interface} and
\ref{prop:body-signed-center-interface}.

\begin{figure}[H]
\centering
\begin{minipage}[t]{.485\textwidth}
\centering
\resizebox{\linewidth}{!}{\begin{tikzpicture}[x=1cm,y=1cm,>=Latex]
  \node[panellabel,anchor=west] at (.05,1.72)
    {(a) selected and actual row coordinates};
  % The two boundary rays meet at the actual regular-hexagon angle \(120^\circ\).
  \coordinate (V) at (5.30,0);
  \coordinate (L) at (4.48,1.420);
  \coordinate (R) at (4.48,-1.420);
  \coordinate (Orow) at (2.30,0);
  \coordinate (Amark) at ($(V)!.42!(L)$);
  \coordinate (Aactual) at ($(V)!.55!(L)$);
  \coordinate (Bmark) at ($(V)!.38!(R)$);
  \coordinate (Ymark) at ($(V)!.64!(R)$);
  \coordinate (Cmark) at ($(V)!.52!(Orow)$);

  \draw[black!45,line width=.65pt] (L)--(V)--(R);
  \draw[ray] (V)--(Orow);
  \fill (V) circle (1pt);
  \node[figlabel,right] at (V) {$V_i$};
  \node[figlabel,above left] at (L) {$V_{i-1}$};
  \node[figlabel,below left] at (R) {$V_{i+1}$};

  \draw[actualreach] (V)--(Aactual);
  \draw[demand] (V)--(Amark);
  \draw[dimline] ($(V)+(0,.18)$)--($(Amark)+(0,.18)$);
  \draw[dimline] ($(Amark)+(0,.18)$)--($(L)+(0,.18)$);
  \node[figlabel,atlasbox,align=left,anchor=east] (topnote) at (4.00,.92)
    {$a\le A$: selected versus actual\\
     $x=1-a$: complementary coordinate};
  \draw[-{Latex[length=1.2mm]},black!45,line width=.45pt]
    (topnote.east)--($(Amark)+(0,.18)$);

  \draw[actualreach] (V)--(Bmark)
    node[midway,below right,figlabel] {$B$};
  \draw[boundcurve,dashed] ($(V)+(0,-.08)$)--($(Ymark)+(0,-.08)$);
  \node[figlabel,atlasbox,anchor=east] (ynote) at (3.95,-.48)
    {$B\le y=g_c(x)$};
  \draw[-{Latex[length=1.2mm]},black!45,line width=.45pt]
    (ynote.east)--($(Bmark)!.58!(Ymark)$);
  \draw[actualreach] (R)--(Ymark);
  \node[figlabel,atlasbox,anchor=east,align=right] (nextnote) at (3.95,-1.13)
    {$1-y=g_c^\vee(a)$\\[-1pt]$\le A_{\rm next}$\\[-1pt]
     center-free outgoing edge};
  \draw[-{Latex[length=1.2mm]},black!45,line width=.45pt]
    (nextnote.east)--($(R)!.48!(Ymark)$);

  \draw[radialdemand] (V)--(Cmark);
  \draw[actualreach] (Cmark)--($(V)!.70!(Orow)$);
  \node[figlabel,above] at ($(V)!.45!(Orow)$) {$c\le C$};
\end{tikzpicture}
}
\end{minipage}\hfill
\begin{minipage}[t]{.485\textwidth}
\centering
\resizebox{\linewidth}{!}{\begin{tikzpicture}[x=1cm,y=1cm,>=Latex]
  \node[panellabel,anchor=west] at (.62,2.45)
    {(b) raw, hatted, and complement-dual views};

  \begin{scope}[shift={(.18,-.02)},x=2.15cm,y=2.15cm]
    \draw[->,black!55] (-.03,0)--(1.08,0) node[right,figlabel] {$x$};
    \draw[->,black!55] (0,-.03)--(0,1.08) node[above,figlabel] {$y$};
    \draw[cappedcurve] (0,0)--(1,1);
    \draw[boundcurve,domain=.015:.985,samples=70,smooth,variable=\x]
      plot ({\x},{(\x-1+sqrt(1+6*\x-3*\x*\x))/2});
    \node[figlabel,text=DemandRed] at (.50,.80) {$g_0(x)>x$};
    \node[figlabel,text=RadialTeal,rotate=38] at (.70,.60)
      {$\widehat g_0=\mathrm I$};
    \node[figlabel,below left] at (0,0) {$0$};
    \node[figlabel,below] at (1,0) {$1$};
    \node[figlabel,left] at (0,1) {$1$};
  \end{scope}

  \node[figlabel,atlasbox] (a) at (3.12,1.22) {$a$};
  \node[figlabel,atlasbox] (x) at (4.18,1.22) {$1-a$};
  \node[figlabel,atlasbox] (fx) at (5.34,1.22) {$f(1-a)$};
  \node[figlabel,atlasbox,align=center] (dual) at (5.34,.22)
    {$f^\vee(a)$\\$=1-f(1-a)$};
  \draw[flow] (a)--node[above,figlabel] {$1-\cdot$}(x);
  \draw[flow] (x)--node[above,figlabel] {$f$}(fx);
  \draw[flow] (fx)--node[left,figlabel] {$1-\cdot$}(dual);
  \node[figlabel] at (4.05,-.35)
    {graph half-turn: $(x,y)\mapsto(1-x,1-y)$};
\end{tikzpicture}
}
\end{minipage}

\medskip
\resizebox{.985\linewidth}{!}{\begin{tikzpicture}[x=1cm,y=1cm,>=Latex]
  \node[panellabel,anchor=west] at (.02,2.30)
    {(c) the center-assisted residual on a skeleton edge};

  \foreach \sx/\pt/\et/\case/\formula in {
    1.85/.25/.25/{p<L}/{1-p},
    6.00/.58/.72/{L\le p<U}/{1-U},
    10.15/.86/.86/{p\ge U}/{1-p}
  }{
    \begin{scope}[shift={(\sx,.65)},scale=1.08]
      \HexCoords
      \DrawHexagon
      \DrawRays
      \coordinate (Pp) at ($(V0)!\pt!(V1)$);
      \coordinate (PL) at ($(V0)!.42!(V1)$);
      \coordinate (PU) at ($(V0)!.72!(V1)$);
      \coordinate (Pe) at ($(V0)!\et!(V1)$);

      \draw[black!28,line width=2.8pt] (V0)--(Pe);
      \draw[demand] (V0)--(Pp);
      \draw[centertrace] (PL)--(PU);
      \draw[radialdemand]
        ($(V1)+(.05,.03)$)--($(Pe)+(.05,.03)$);
      \fill[DemandRed] (Pp) circle (.75pt);
      \fill[black!65] (Pe) circle (.7pt);

      \node[figlabel,above right] at ($(V0)!.52!(Pp)$) {$p$};
      \node[figlabel,above] at ($(PL)!.50!(PU)$) {$J=[L,U]$};
      \node[figlabel,anchor=east,align=right] at ($(V1)+(-.06,.25)$)
        {$A_{\rm next}\ge\mathcal R_J(p)$};
      \node[figlabel,below] at (0,-1.55)
        {$\case:\quad\mathcal R_J(p)=\formula$};
    \end{scope}
  }

  \node[figlabel,align=center] at (6.00,-1.78)
    {Gray is the component of \([0,p]\cup J\) attached to \(V_0\);
     the teal far-side reach is forced from \(V_1\).\\
     For \(p=g_c(1-a)\) this gives \(g_{c,J}^{\vee}(a)\);
     for \(p=\widehat g_c(1-a)\) it gives
     \(\widehat g_{c,J}^{\vee}(a)\).};
\end{tikzpicture}
}

\medskip
\begin{minipage}[t]{.485\textwidth}
\centering
\resizebox{\linewidth}{!}{\begin{tikzpicture}[x=1cm,y=1cm,>=Latex]
  \node[panellabel,anchor=west] at (1.55,2.35)
    {(d) strict-supercritical envelope};
  \begin{scope}[shift={(.22,-.02)},x=2.12cm,y=2.12cm]
    \draw[->,black!55] (-.03,0)--(1.08,0) node[right,figlabel] {$x$};
    \draw[->,black!55] (0,-.03)--(0,1.08) node[above,figlabel] {$g_c(x)$};
    \draw[inactive] (0,0)--(1,1);
    \draw[boundcurve]
      (0,0)--(.22,.22)
      ..controls(.32,.34) and (.40,.72)..(.64,.83)
      ..controls(.72,.86) and (.80,.85)..(.843,.843);
    \draw[inactive] (0,.843)--(1,.843);
    \fill[white,draw=DemandRed,line width=.9pt] (.843,.843) circle (.018);
    \node[figlabel,text=GapPurple] at (.49,.53) {$g_c(x)>x$};
    \node[figlabel,anchor=west] at (.07,.90) {$g_c^{\rm sc}$ (supremum)};
    \node[figlabel,below] at (.843,0) {$x_*$};
  \end{scope}

  \node[figlabel,atlasbox,align=center] (sc) at (3.22,1.16)
    {strict row\\$B<g_c^{\rm sc}$};
  \node[figlabel,atlasbox,align=center] (next) at (5.25,1.16)
    {next row\\$A_{\rm next}>1-g_c^{\rm sc}$};
  \draw[flow] (sc)--(next);
  \node[figlabel] at (4.25,1.79) {center-free edge};
  \draw[centertrace] (3.55,.28)--(4.35,.28);
  \node[figlabel,above] at (3.95,.30) {$J$};
  \draw[DemandRed,line width=.9pt] (3.78,.11)--(4.12,.45);
  \draw[DemandRed,line width=.9pt] (3.78,.45)--(4.12,.11);
  \node[figlabel,align=left,anchor=west] at (4.50,.28)
    {with a center trace,\\the complement need not follow};
\end{tikzpicture}
}
\end{minipage}\hfill
\begin{minipage}[t]{.485\textwidth}
\centering
\resizebox{\linewidth}{!}{\begin{tikzpicture}[x=1cm,y=1cm,>=Latex]
  \node[panellabel,anchor=west] at (.03,1.72)
    {(e) slotwise lower relaxation};
  \foreach \x/\i in {.55/0,2.15/1,3.75/2,5.35/3}{
    \node[figlabel,atlasbox] (x\i) at (\x,1.10) {$x_{\i}$};
    \node[figlabel,atlasbox] (y\i) at (\x,.05) {$y_{\i}$};
  }
  \draw[inactive] (y0)--node[right,figlabel] {$=$}(x0);
  \foreach \i in {1,2,3}{
    \draw[inactive,->] (y\i)--node[right,figlabel] {$\le$}(x\i);
  }
  \foreach \i/\j in {0/1,1/2,2/3}{
    \draw[flow] (x\i)--node[above,figlabel] {$\Phi_{\j}$}(x\j);
    \draw[cappedcurve,->] (y\i)--node[below,figlabel]
      {$\underline\Phi_{\j}$}(y\j);
  }
  \node[figlabel,align=center] at (2.95,-.62)
    {$y_3=[\underline\Phi_1\mid\underline\Phi_2
    \mid\underline\Phi_3](x_0)\le x_3$;\quad leftmost slot acts first.};
\end{tikzpicture}
}
\end{minipage}
\caption{Reading the transfer alphabet.  Panel (a) distinguishes selected
demands from actual reaches: $x=1-a$ is a coordinate defect, and $g_c(x)$ is
an extremal admissible output rather than the actual $B$.  Panel (b) is the
exact zero-radial graph and shows that the identity is
$\widehat g_0$, not raw $g_0$; the right-hand diagram explains $\vee$ by
complementing both coordinates.  Panel (c) shows that only the component
attached to $0$ enters $\mathcal R_J$.  Panel (d) is schematic: the
strict-supercritical supremum is not attained, and its complementary
next-reach bound needs a center-free edge.  Panel (e) shows where
nondecreasing lower relaxations enter a composition.}
\label{fig:transfer-notation-atlas}
\end{figure}

\clearpage

\begin{figure}[H]
\centering
\begin{minipage}[t]{.485\textwidth}
\centering
\resizebox{\linewidth}{!}{\begin{tikzpicture}[x=1cm,y=1cm,>=Latex]
  \node[panellabel,anchor=west] at (.03,2.12)
    {(f) signed parameter cascade};
  \node[figlabel,atlasbox] (R) at (.28,1.25) {$R$};
  \node[figlabel,atlasbox] (W) at (1.55,1.25) {$W=1-R$};
  \node[figlabel,atlasbox] (E) at (3.62,1.25) {$E=\sqrt{1-RW}$};
  \node[figlabel,atlasbox] (eta) at (5.55,1.67) {$\eta=1-E$};
  \node[figlabel,atlasbox] (P) at (5.55,.84) {$P=E\eta$};
  \draw[flow] (R)--(W);
  \draw[flow] (W.east)--(E.west);
  \draw[flow] (E)--(eta);
  \draw[flow] (E)--(P);

  \node[figlabel,atlasbox,align=center] (slack) at (1.50,.02)
    {$\alpha=F_0(O),\ \delta=F_2(O)$\\
     $k=\eta+\alpha+\delta$};
  \node[figlabel,atlasbox,align=center] (Delta) at (4.78,-.02)
    {$\Delta_R=P-\alpha-W\delta$\\
     $\Delta_L=P-R\alpha-\delta$};
  \draw[flow] (slack)--(Delta);
  \draw[flow] (P)--(Delta);
\end{tikzpicture}
}
\end{minipage}\hfill
\begin{minipage}[t]{.485\textwidth}
\centering
\resizebox{\linewidth}{!}{\begin{tikzpicture}[x=1cm,y=1cm,>=Latex]
  \node[panellabel,anchor=west] at (.03,2.12)
    {(g) the sign of the companion surplus};
  \foreach \yy/\leftpt/\rightpt/\name/\desc in {
    1.15/3.15/2.35/{$\Delta_L<0$}/{reversed: CE1, no trace},
    .43/2.75/2.75/{$\Delta_L=0$}/{point contact: CE1},
    -.29/2.10/3.42/{$\Delta_L>0$}/{positive trace: CE2}
  }{
    \draw[black!50,line width=.6pt] (.55,\yy)--(4.75,\yy);
    \draw[inactive] (\leftpt,\yy)--(\rightpt,\yy);
    \fill[black!60] (\leftpt,\yy) circle (.9pt);
    \fill[black!60] (\rightpt,\yy) circle (.9pt);
    \node[figlabel,left] at (.50,\yy) {\name};
    \node[figlabel,right,align=left] at (4.78,\yy) {\desc};
  }
  \draw[centertrace] (2.10,-.29)--(3.42,-.29);
  \node[figlabel,above] at (2.76,-.27)
    {$\ell_L=[\Delta_L]_+/W$};
  \node[figlabel,below] at (2.35,1.13) {$R+\alpha$};
  \node[figlabel,above] at (3.15,1.17) {$k/W$};
\end{tikzpicture}
}
\end{minipage}

\medskip
\resizebox{.985\linewidth}{!}{\begin{tikzpicture}[x=1cm,y=1cm,>=Latex]
  \node[panellabel,anchor=west] at (.02,2.16)
    {(h) two exact signed-center examples on the full skeleton};

  \begin{scope}[shift={(2.75,.30)},scale=1.42]
    \HexCoords
    \DrawHexagon
    \DrawRays
    \draw[centerrole]
      (.679270858,-.541343219)--(.793978725,.452056048)--
      (-.123684211,.054696341)--cycle;
    \fill[CenterBlue] (O) circle (1pt);

    \draw[centertrace]
      (.784816491,.372708771)--(.75,.433012702);
    \draw[inactive]
      (.677224736,-.559063157)--(.685,-.545596004);
    \fill[black!55] (.677224736,-.559063157) circle (.75pt);
    \fill[black!55] (.685,-.545596004) circle (.75pt);

    \node[figlabel,above right] at (.765,.414) {$\ell_R$};
    \node[figlabel,anchor=east,align=right] at (.53,-.56)
      {reversed candidate\\endpoints};
    \node[figlabel,right] at (V0) {$V_0$};
    \node[figlabel,above] at ($(V0)!.55!(V1)$) {$e_{0,1}$};
    \node[figlabel,below right] at ($(V0)!.55!(V5)$) {$e_{5,0}$};
  \end{scope}
  \node[figlabel,align=center] at (2.75,-1.47)
    {\(\mathrm{CE1}\): \(R=\frac35,\ W=\frac25\),\\
     \(\alpha=F_0(1,1)=\frac3{100},\quad
       \delta=F_2(1,1)=\frac1{10}\),\\
     \(\Delta_R>0>\Delta_L,\quad
       \ell_R\approx.06963,\quad\ell_L=0\).};

  \begin{scope}[shift={(8.70,.30)},scale=1.42]
    \HexCoords
    \DrawHexagon
    \DrawRays
    \draw[centerrole]
      (.734534016,-.582365475)--(.849241883,.411033792)--
      (-.068421053,.013674085)--cycle;
    \fill[CenterBlue] (O) circle (1pt);

    \draw[centertrace]
      (.834816491,.286106231)--(.78,.381051178);
    \draw[centertrace]
      (.752224736,-.429159346)--(.685,-.545596004);

    \node[figlabel,above right] at (.80,.35) {$\ell_R$};
    \node[figlabel,below right] at (.71,-.50) {$\ell_L$};
    \node[figlabel,right] at (V0) {$V_0$};
    \node[figlabel,above] at ($(V0)!.55!(V1)$) {$e_{0,1}$};
    \node[figlabel,below right] at ($(V0)!.55!(V5)$) {$e_{5,0}$};
  \end{scope}
  \node[figlabel,align=center] at (8.70,-1.47)
    {\(\mathrm{CE2}\): \(R=\frac35,\ W=\frac25\),\\
     \(\alpha=F_0(1,1)=\frac3{100},\quad
       \delta=F_2(1,1)=\frac1{25}\),\\
     \(\Delta_R,\Delta_L>0,\quad
       \ell_R\approx.10963,\quad\ell_L\approx.13445\).};

  \node[figlabel,align=center] at (5.73,-2.25)
    {\(\alpha,\delta\) are affine side-slack values at \(O\), not drawn lengths;
     the blue boundary segments are the metric traces.};
\end{tikzpicture}
}

\medskip
\begin{minipage}[t]{.485\textwidth}
\centering
\resizebox{\linewidth}{!}{\begin{tikzpicture}[x=1cm,y=1cm,>=Latex]
  \node[panellabel,anchor=west] at (.03,2.03)
    {(i) six exits, measured from $O$};
  \begin{scope}[shift={(1.64,.50)},scale=1.35]
    \HexCoords
    \DrawHexagon
    \DrawRays
    \draw[centerrole]
      (.734534016,-.582365475)--(.849241883,.411033792)--
      (-.068421053,.013674085)--cycle;
    \foreach \i/\d in {
      0/.801779789,1/.066666667,2/.04,3/.05,4/.03,5/.075
    }{
      \coordinate (D\i) at ($(O)!\d!(V\i)$);
      \draw[radialdemand] (O)--(D\i);
      \fill[RadialTeal] (D\i) circle (.9pt);
    }
    \fill[CenterBlue] (O) circle (1pt);
    \node[figlabel,above] at ($(O)!.52!(D0)$) {$d_0^C$};
    \node[figlabel,above right] at ($(O)!.55!(V1)$) {$d_1^C$};
    \node[figlabel,above left] at ($(O)!.55!(V2)$) {$d_2^C$};
    \node[figlabel,left] at ($(O)!.60!(V3)$) {$d_3^C$};
    \node[figlabel,below left] at ($(O)!.55!(V4)$) {$d_4^C$};
    \node[figlabel,below right] at ($(O)!.55!(V5)$) {$d_5^C$};
  \end{scope}
  \node[figlabel,anchor=west,align=left] at (3.25,1.24)
    {$d_0^C=E-\alpha-\delta$\\
     $d_1^C=\delta/R$\\
     $d_2^C=\delta$};
  \node[figlabel,anchor=west,align=left] at (3.25,.05)
    {$d_3^C=\min\{\alpha/R,\delta/W\}$\\
     $d_4^C=\alpha$\\
     $d_5^C=\alpha/W$};
  \node[figlabel,align=center] at (2.72,-1.02)
    {CE2 sample: $(d_0^C,\ldots,d_5^C)
     \approx(.8018,.0667,.04,.05,.03,.075)$};
\end{tikzpicture}
}
\end{minipage}\hfill
\begin{minipage}[t]{.485\textwidth}
\centering
\resizebox{\linewidth}{!}{\begin{tikzpicture}[x=1cm,y=1cm,>=Latex]
  \node[panellabel,anchor=west] at (.03,2.03)
    {(j) scalar boundary-contribution bounds};

  \node[figlabel,anchor=east] at (1.28,1.19) {$\mathrm{CE1}$};
  \draw[black!50,line width=.6pt] (1.45,1.19)--(5.35,1.19);
  \draw[demand] (1.45,1.19)--(3.78,1.19);
  \draw[black!60] (3.78,1.08)--(3.78,1.30);
  \draw[black!60] (4.73,1.08)--(4.73,1.30);
  \draw[black!60] (5.18,1.08)--(5.18,1.30);
  \node[figlabel,below] at (3.78,1.08) {$L_{\partial H}=\ell_R$};
  \node[figlabel,above] at (4.73,1.30) {$P$};
  \node[figlabel,above] at (5.18,1.30)
    {$\sqrt3/2-3/4$};
  \node[figlabel,align=center] at (3.40,.62)
    {$L_{\partial H}(T_C)\le P\le\sqrt3/2-3/4$};

  \node[figlabel,anchor=east] at (1.28,-.05) {$\mathrm{CE2}$};
  \draw[black!50,line width=.6pt] (1.45,-.05)--(5.35,-.05);
  \draw[demand] (1.45,-.05)--(2.30,-.05);
  \draw[centertrace] (2.30,-.05)--(3.34,-.05);
  \draw[black!60] (5.18,-.16)--(5.18,.06);
  \node[figlabel,above] at (1.88,.02) {$\ell_R$};
  \node[figlabel,above] at (2.82,.02) {$\ell_L$};
  \node[figlabel,below] at (3.34,-.14) {$L_{\partial H}=\ell_R+\ell_L$};
  \node[figlabel,above] at (5.18,.06) {$1/2$};
  \node[figlabel,align=center] at (3.40,-.72)
    {$L_{\partial H}(T_C)<1/2$};
\end{tikzpicture}
}
\end{minipage}
\caption{The signed CE1/CE2 interface.  Panel (f) records which parameters are
derived and identifies $\alpha,\delta$ as affine side-slack values at $O$.
Panel (g) explains the positive part $[\Delta_L]_+$ and keeps the
$\Delta_L=0$ point-contact state in CE1.  Panel (h) uses the valid samples
$R=3/5$, $W=2/5$, $\alpha=3/100$, with $\delta=1/10$ for CE1 and
$\delta=1/25$ for CE2; rounded coordinates are used only to draw the exact
unit triangles.  Panel (i) maps each exit $d_i^C$ to its ray and measures it
from $O$.  Panel (j) consists of scalar comparison bars: they are not
locations on $\partial H$, and $L_{\partial H}(T_C)$ is not the perimeter of
$T_C$.}
\label{fig:signed-center-notation-atlas}
\end{figure}
